\documentclass[10pt,a4paper]{amsart}
\usepackage[margin=19mm]{geometry}
\usepackage{amsmath,amssymb,mathtools}
\usepackage[scr=rsfs,cal=euler]{mathalfa}
\usepackage{xfrac}
\usepackage[unicode,hidelinks,bookmarks=false]{hyperref}
\newcommand{\ie}{i.e.\ }
\newcommand{\cf}{cf.\ }
\newcommand{\resp}{respectively\ }
\newcommand{\Subsection}{\subsection}
\providecommand{\nobreakdash}{\nobreak\hbox{-}}
\setlength{\emergencystretch}{2em}
\multlinegap0pt
\usepackage{bm}
\usepackage{amssymb}
\usepackage{float}
\usepackage{xcolor}
\newtheorem{proposition}{Proposition}
\title{A Heuristic Alternating Direction Method of Multipliers Framework for Distributed and Centralized Tree-Constrained Optimization:\\ Applications to Hop-Constrained Spanning Tree Multicommodity Flow Design}
\author{Yacine Mokhtari}
\date{Source: arXiv:2508.11078v2 (CC BY 4.0)}
\begin{document}
\maketitle

\noindent\emph{Source and licence.} A Heuristic Alternating Direction Method of Multipliers Framework for Distributed and Centralized Tree-Constrained Optimization:\\ Applications to Hop-Constrained Spanning Tree Multicommodity Flow Design. Version arXiv:2508.11078v2 (CC BY 4.0), distributed by arXiv under the Creative Commons Attribution 4.0 International licence, http://creativecommons.org/licenses/by/4.0/, as recorded in the arXiv licence metadata for this exact version and shown on the abstract page https://arxiv.org/abs/2508.11078v2. Full licence notice: \texttt{licenses/arxiv-2508.11078-CC-BY-4.0.txt}.\par\smallskip

\noindent\textit{Editorial scope.} Counted unit: the article's proposition Proposition (source0.tex lines 221--232). The article's complete original proof of it is delivered with it (source0.tex lines 234--268, one \texttt{proof} environment of 35 lines). No mathematics is written, repaired or re-derived: the statement and the proof are the article's own lines, in the article's own notation, and no retained line is substituted or re-flowed.  Retained uncounted as the source-local context the proof needs: lines 206--215.  Nothing was omitted from any retained span, so the package carries no omission record.  No retained line contains a citation, so the wrapper carries no bibliography and no bibliography entry.  No retained line contains a figure environment, an image-inclusion command or drawing code, so no image is delivered and the package declares no asset dependency.  Editorial formatting only, in addition: an AMS \texttt{amsart} preamble that restates the article's own declarations and macros used by the retained lines, namely the environments \texttt{proposition} on separate counters, together with the standard packages those lines need.  The retained environments print in this document as Proposition 1 in order of appearance, whereas the article prints the same items with the numbering of its own full document.  The complete native source of the article as distributed in the arXiv e-print for this exact version is bundled unchanged as \texttt{sources/183295/source0.tex}.
		\begin{align}
			(\boldsymbol{x},\boldsymbol{w})_{k+1}& =\underset{(\boldsymbol{x},%
				\boldsymbol{w})\in \Sigma }{\arg \min }\left[ f(\boldsymbol{x},\boldsymbol{w}%
			)+\frac{\rho }{2}\Vert \boldsymbol{z}_{k}-\boldsymbol{w}+\boldsymbol{\mu }%
			_{k}\Vert^{2}\right] ,  \label{AA1} \\[1ex]
			\boldsymbol{z}_{k+1}& =\underset{\boldsymbol{z}\in \mathcal{Z}}{\arg \min }%
			\Vert \boldsymbol{z}-\boldsymbol{w}_{k+1}+\boldsymbol{\mu }_{k}\Vert^{2},  \label{AA2} \\[1ex]
			\boldsymbol{\mu }_{k+1}& =\boldsymbol{\mu }_{k}+\boldsymbol{z}_{k+1}-%
			\boldsymbol{w}_{k+1}.  \label{AA3}
		\end{align}

	\begin{proposition}
		\label{Proposition} Problem~\eqref{AA2} is equivalent to solving the
		following discrete optimization problem at each iteration: 
		\begin{equation}
			\boldsymbol{z}_{k+1}=\underset{\boldsymbol{z}\in \mathcal{Z}}{\arg \min }\;%
			\boldsymbol{z}^{T}\boldsymbol{h}_{k},  \label{weight}
		\end{equation}%
		where the weight vector $\boldsymbol{h}_{k}\in \mathbb{R}^{m}$ is given by 
		\begin{equation}
			\boldsymbol{h}_{k}=\boldsymbol{\mu }_{k}-\boldsymbol{w}_{k+1}.  \label{h_k}
		\end{equation}%
	\end{proposition}

	\begin{proof}
		\label{Proof_prop} We begin by expanding the squared norm appearing in the
		second ADMM subproblem: We aim to show that problem~\eqref{AA2} is
		equivalent to solving the discrete optimization problem~\eqref{weight}.
		Starting with the objective in~\eqref{AA2}, we expand the squared norm: 
		\begin{equation*}
			\left\Vert \boldsymbol{z}-\boldsymbol{w}_{k+1}+\boldsymbol{\mu }%
			_{k}\right\Vert ^{2}=\boldsymbol{z}^{\top }\boldsymbol{z}+\left\Vert 
			\boldsymbol{\mu }_{k}-\boldsymbol{w}_{k+1}\right\Vert ^{2}+2\boldsymbol{z}%
			^{\top }\left( \boldsymbol{\mu }_{k}-\boldsymbol{w}_{k+1}\right) .
		\end{equation*}%
		Since $\boldsymbol{z}\in \mathcal{Z}\subset \{0,1\}^{m}$ represents a
		spanning tree or arborescence with exactly $n-1$ edges, we have $\boldsymbol{%
			z}^{\top }\boldsymbol{z}=\sum_{(i,j)\in \mathcal{E}\text{ or }\mathcal{A}%
		}z_{ij}=n-1$, as each selected edge contributes 1 to the sum. Thus, the
		expression becomes: 
		\begin{eqnarray*}
			\left\Vert \boldsymbol{z}-\boldsymbol{w}_{k+1}+\boldsymbol{\mu }%
			_{k}\right\Vert ^{2} &=&(n-1)+\left\Vert \boldsymbol{\mu }_{k}-\boldsymbol{w}%
			_{k+1}\right\Vert ^{2}+2\boldsymbol{z}^{\top }\left( \boldsymbol{\mu }_{k}-%
			\boldsymbol{w}_{k+1}\right)  \\
			&=&(n-1)+\left\Vert \boldsymbol{\mu }_{k}-\boldsymbol{w}_{k+1}\right\Vert
			^{2}+2\boldsymbol{z}^{\top }\boldsymbol{h}_{k}.
		\end{eqnarray*}%
		Since $n-1$ and $\left\Vert \boldsymbol{\mu }_{k}-\boldsymbol{w}%
		_{k+1}\right\Vert ^{2}$ are constant with respect to $\boldsymbol{z}$,
		minimizing the expression over $\boldsymbol{z}\in \mathcal{Z}$ is equivalent
		to minimizing $\boldsymbol{z}^{\top }\boldsymbol{h}_{k}$. Thus, problem~%
		\eqref{AA2} reduces to: 
		\begin{equation*}
			\boldsymbol{z}_{k+1}=\underset{\boldsymbol{z}\in \mathcal{Z}}{\arg \min }\ 
			\boldsymbol{z}^{\top }\boldsymbol{h}_{k},
		\end{equation*}%
		which matches~\eqref{weight}.
	\end{proof}

\end{document}
